%% kapa: exam=2 round=30 year=2019 period=2 subject=appraisal_theory subject_ko=감정평가이론 seq=3 number=3 points=20
%% source: Q-Net 감정평가사 기출문제 (cst00309) artlSeq=5205595 file=raw/1008/5205595/감정평가이론.pdf
%% conversion: kordoc 4.24.51 -> tools/gen2.py (auto); review=proofread-gpt-terra; flags=none
\begin{longquestion}{【문제 3】}
「감정평가에 관한 규칙」에서 감정평가 시 시장가치기준을 원칙으로 하되, 예외적인 경우 ‘시장가치 외의 가치’를 인정하고 있다. 그러나 현행 「감정평가에 관한 규칙」에서는 ‘시장가치 외의 가치’에 대한 유형 등의 구체적인 설명이 없어 이를 보완할 필요성이 있다. 감정평가 시 적용할 수 있는 구체적인 ‘시장가치 외의 가치’에 대해 설명하시오. (20점) \par
\end{longquestion}
